\section{DiffForm degree-probe alignment}
\label{sec:diffform-degree-probe-alignment}

\begin{definition}[Differential-form degree-probe alignment row]
\label{def:diffform-degree-probe-alignment-row}
For a unary degree history $d$ and a visible probe bundle $\mathcal{P}$, the
row $\mathsf{DegreeProbeAligned}(d,\mathcal{P})$ is the closed structural
relation generated by the following constructors:
\begin{enumerate}
  \item $\mathsf{DegreeProbeAligned}(\emp,\mathsf{ProbeBundle.Bnil})$.
  \item If $\mathsf{DegreeProbeAligned}(d,\mathcal{P})$ holds, and if $p$ is a
  tangent-probe row admitted by the manifold source, then
  $\mathsf{DegreeProbeAligned}(\Eone(d),\mathsf{ProbeBundle.Bcons}(p,\mathcal{P}))$.
\end{enumerate}
Thus the unary degree row and the $\mathsf{ProbeBundle}$ spine expose the same
visible length. The relation is part of the displayed differential-form row:
it does not import a host tangent-space dimension and does not allocate a
second arity counter outside $BHist$.
\end{definition}
\leandef{BEDC.Derived.DiffFormUp.DegreeProbeAligned}

\begin{theorem}[Differential-form carrier refines to degree-probe alignment]
\label{thm:diffform-degree-probe-alignment-carrier-refinement}
Let
$$
\omega=(d,\mathcal{P},\tau,\sigma,\alpha,\lambda)
$$
be carried by $\mathsf{DiffFormBHistCarrier}_{M,T}$. Suppose that the tensor
row $\tau$ is read over exactly the displayed bundled tangent probes in
$\mathcal{P}$, with one unary degree constructor in $d$ for each visible probe
slot and with no tensor input outside the bundle. Then $\omega$ supplies
$\mathsf{DegreeProbeAligned}(d,\mathcal{P})$.
\end{theorem}
\begin{proof}
First unpack the carrier through
\autoref{def:diffform-bhist-form-carrier}. The carrier exposes
$\mathsf{UnaryHistory}(d)$, the finite $\mathsf{ProbeBundle}$ spine
$\mathcal{P}$, and the tensor row $\tau$ over that same bundle.

Induct simultaneously on the visible probe spine and on the unary degree row
selected by the tensor-input reading. In the empty bundle case, the carrier
clauses and the exact-reading hypothesis leave no tensor input slot. The
unary row is therefore the empty degree row, and the empty constructor of
\autoref{def:diffform-degree-probe-alignment-row} applies.

In the cons case, the carrier exposes a head tangent probe and a tail
$\mathsf{ProbeBundle}$. The exact-reading hypothesis strips one $\Eone$
constructor from the degree row for that head probe and leaves the tail tensor
row over the tail bundle. The induction hypothesis gives alignment for the
tail, and the $\mathsf{Bcons}$ constructor of
\autoref{def:diffform-degree-probe-alignment-row} rebuilds alignment for the
whole row.

Finally apply
\autoref{thm:diffform-degree-probe-carrier-admissibility}. It says that the
carried row supplies the displayed degree row and the visible probe spine as
the support surface consumed by the tensor, scalar, antisymmetry, and source
ledgers. Hence no hidden probe slot or hidden degree constructor remains
outside the aligned row.
\end{proof}
\leanchecked{BEDC.Derived.DiffFormUp.DiffFormDegreeProbeAligned\_empty\_iff}

\begin{theorem}[Differential-form degree-probe alignment transport]
\label{thm:diffform-degree-probe-alignment-transport}
Let
$$
\omega=(d,\mathcal{P},\tau,\sigma,\alpha,\lambda)
$$
and
$$
\omega'=(d',\mathcal{P}',\tau',\sigma',\alpha',\lambda')
$$
be rows carried by $\mathsf{DiffFormBHistCarrier}_{M,T}$ with
$\omega\sim_{\DiffFormUp}\omega'$. If
$\mathsf{DegreeProbeAligned}(d,\mathcal{P})$ holds, then
$\mathsf{DegreeProbeAligned}(d',\mathcal{P}')$ holds.
\end{theorem}
\leanchecked{BEDC.Derived.DiffFormUp.DiffFormDegreeProbeAligned\_hsame\_transport}
\begin{proof}
Project the classifier comparison through
\autoref{def:diffform-bhist-form-classifier}. The degree clause gives
$\hsame(d,d')$, while the probe and tensor clauses match the displayed
$\mathsf{ProbeBundle}$ spines componentwise through the manifold and tensor
sources.

Apply \autoref{thm:diffform-degree-probe-transport}. The support row carried
by $\omega$ transports to the support row carried by $\omega'$, with the same
visible probe-slot arity.

Now apply \autoref{thm:diffform-probe-slot-arity-transport}. Since the probe
slot count is preserved and the unary degree row is transported by $\hsame$,
each constructor step in
\autoref{def:diffform-degree-probe-alignment-row} has a matching transported
constructor step. Rebuilding the transported spine gives
$\mathsf{DegreeProbeAligned}(d',\mathcal{P}')$.
\end{proof}

\begin{theorem}[Differential-form wedge preserves degree-probe alignment]
\label{thm:diffform-degree-probe-alignment-wedge-closure}
Let $\omega$ and $\eta$ be carried $\DiffFormUp$ rows with
$\mathsf{DegreeProbeAligned}(d_\omega,\mathcal{P}_\omega)$ and
$\mathsf{DegreeProbeAligned}(d_\eta,\mathcal{P}_\eta)$. If the wedge row
$\omega\wedge_{\DiffFormUp}\eta$ is admitted by
$\mathsf{DiffFormBHistCarrier}_{M,T}$ with output degree
$d_{\omega\wedge\eta}$ and output probe bundle
$\mathcal{P}_{\omega\wedge\eta}$, then
$\mathsf{DegreeProbeAligned}(d_{\omega\wedge\eta},
\mathcal{P}_{\omega\wedge\eta})$ holds.
\end{theorem}
\begin{proof}
Expose the admitted wedge probe surface with
\autoref{thm:diffform-degree-probe-wedge-compatibility}. The output bundle is
the visible append
$$
\mathcal{P}_{\omega\wedge\eta}
=\mathsf{bundleAppend}(\mathcal{P}_\omega,\mathcal{P}_\eta),
$$
and the tensor row consumes exactly that appended spine.

Induct on the left probe bundle $\mathcal{P}_\omega$. In the empty case, the
left alignment row identifies $d_\omega$ with the empty unary degree, and the
append surface reduces to the right input probe bundle. The $\Cont$ left-unit
row in \autoref{def:diffform-wedge-degree-ledger} identifies the output degree
with $d_\eta$, so the right input alignment supplies the output alignment.

In the cons case, the left alignment row strips one head probe from
$\mathcal{P}_\omega$ and one $\Eone$ constructor from $d_\omega$. The induction
hypothesis aligns the wedge of the left tail with $\eta$. Reapplying the
$\mathsf{Bcons}$ constructor in
\autoref{def:diffform-degree-probe-alignment-row} aligns the restored head
probe with the restored unary constructor.

It remains to justify that the degree endpoint is the wedge endpoint rather
than an auxiliary row. Use the continuation row
$$
\Cont(d_\omega,d_\eta,d_{\omega\wedge\eta})
$$
from \autoref{def:diffform-wedge-degree-ledger}. Continuation associativity
and the unit rows already cited by that ledger identify the inductive endpoint
with the admitted wedge degree. Repacking with the visible appended bundle
gives $\mathsf{DegreeProbeAligned}(d_{\omega\wedge\eta},
\mathcal{P}_{\omega\wedge\eta})$.
\end{proof}
\leanchecked{BEDC.Derived.DiffFormUp.DiffFormDegreeProbeAligned\_bundleAppend\_cont}
